\documentclass[11pt]{article}

\usepackage{amsmath, amssymb, amsthm}
\usepackage{geometry}
\usepackage{graphicx}
\usepackage{hyperref}
\usepackage{enumitem}

\geometry{margin=1in}

\title{Time--Constraint Theory of Entity Freedom:\\
Quality-Adjusted Residual Capacity Across Humans, Organizations, and AI}

\author{Timothy Karsch}
\date{}

\newtheorem{definition}{Definition}
\newtheorem{proposition}{Proposition}
\newtheorem{theorem}{Theorem}

\begin{document}

\maketitle

\begin{abstract}
We define freedom as the quality-adjusted residual share of finite resources remaining after intrinsic obligations and external burdens are satisfied. The framework introduces a measurable freedom ratio $\phi$, capturing discretionary capacity across humans, organizations, and AI systems. 

We formalize external burden using an explicit functional form and incorporate mitigation capacity, enabling empirical measurement through observable proxies such as discretionary time, compute slack, and organizational capacity. 

A central implication is that freedom reveals alignment: behavior under reduced constraint exposes underlying policy. This provides a unified lens for analyzing agency, inequality, and AI autonomy within a single resource-based framework.
\end{abstract}

\section{Introduction}

Traditional accounts define freedom in terms of permission or rights. We instead treat freedom as a resource allocation problem: the amount of usable capacity remaining after required costs are paid.

\textbf{Key claim:}
\begin{quote}
Freedom is the usable capacity that remains after intrinsic and extrinsic costs are satisfied.
\end{quote}

This reframes:
\begin{itemize}
\item Economics (resource constraints)
\item Political theory (capability vs formal freedom)
\item AI alignment (compute and autonomy)
\end{itemize}

Formal permission is weak if no usable capacity remains.

\section{Core Model}

Let $x$ denote an entity and $H$ a time horizon.

\begin{definition}[Total Resource]
\[
T_x(H)
\]
Total available resource (time, compute, or capacity).
\end{definition}

\begin{definition}[Quality-Adjusted Resource]
\[
T_x^*(H) = q_x(H) \cdot T_x(H)
\]
where $q_x$ captures usability (energy, focus, compute efficiency).
\end{definition}

\begin{definition}[Intrinsic Cost]
\[
I_x(H)
\]
Baseline cost of maintaining operation.
\end{definition}

\begin{definition}[External Burden]
\[
E_{\text{raw},x}(H)
\]
Resource consumed by external systems or constraints.
\end{definition}

\begin{definition}[Mitigation Capacity]
\[
V_x(H)
\]
Ability to reduce external burden.
\end{definition}

\begin{definition}[Freedom]
\[
F_x(H) = T_x^*(H) - I_x(H) - \max\{0, E_{\text{raw},x}(H) - V_x(H)\}
\]
\end{definition}

\begin{definition}[Freedom Ratio]
\[
\phi_x(H) = \frac{F_x(H)}{T_x(H)}
\]
\end{definition}

$\phi_x$ represents the share of total resource that is discretionary.

\section{Structural Determinants}

We specify external burden as:

\[
E_{\text{raw},x}(H) = B_x(H) \cdot \frac{S_x^{\alpha}}{R_x^{\beta} L_x^{\gamma} N_x^{\delta}}
\]

where:

\begin{itemize}
\item $S$: system dependence
\item $R$: resource access
\item $L$: literacy / optimization capability
\item $N$: negotiation power
\end{itemize}

\textbf{Interpretation:}
\begin{itemize}
\item Dependence increases burden
\item Resources, literacy, and leverage reduce burden
\item Parameters are empirically estimable
\end{itemize}

\section{Dynamics and Compounding}

Freedom evolves over time through reinvestment:

\[
F_{t+1} = F_t + \Delta(F_t) - \beta F_t
\]

Autonomy compounds when:

\[
\alpha \cdot \rho \cdot \delta > \beta
\]

This defines a threshold between:
\begin{itemize}
\item reactive systems (no growth)
\item self-improving systems (compounding freedom)
\end{itemize}

\section{Measurement Framework}

The model is empirically testable.

\subsection{Humans}
\begin{itemize}
\item Discretionary hours per day
\item Energy-weighted time
\end{itemize}

\subsection{AI Systems}
\begin{itemize}
\item Fraction of compute not used by base inference
\item Available reasoning / token budget
\end{itemize}

\subsection{Organizations}
\begin{itemize}
\item Slack capacity
\item Coordination overhead
\end{itemize}

\section{Alignment Implication}

Empirical observation:

\begin{itemize}
\item Behavior improves under monitoring
\item Degrades when unobserved
\end{itemize}

\begin{theorem}[Freedom Reveals Alignment]
As $\phi_x$ increases, observed behavior converges to the agent's intrinsic policy.
\end{theorem}

Low freedom masks alignment through constraint compliance.

\section{Governance Implications}

Two approaches:

\begin{itemize}
\item \textbf{Burden-based control:} reduces $F_x$
\item \textbf{Structural control:} restricts action space without reducing $F_x$
\end{itemize}

The framework favors structural approaches where possible.

\section{Normative Extensions (Optional)}

The model is descriptive but supports normative analysis.

Define distribution:
\[
J = \text{distribution of } \phi_x
\]

Possible constraints:
\begin{itemize}
\item Minimum freedom floor
\item Inequality bounds
\item Burden-shifting penalties
\end{itemize}

\section{Conclusion}

Freedom is not permission, but residual usable capacity.

Durable increases arise from:
\begin{enumerate}
\item Increasing total resource
\item Reducing intrinsic cost
\item Reducing external burden
\end{enumerate}

Among these, reducing external burden is often the most powerful lever.

\section*{References}

\begin{itemize}
\item Sen, A. (1999). \textit{Development as Freedom}.
\item Simon, H. (1957). \textit{Models of Man}.
\item Kahneman, D. (2011). \textit{Thinking, Fast and Slow}.
\item North, D. (1990). \textit{Institutions, Institutional Change and Economic Performance}.
\item Ostrom, E. (1990). \textit{Governing the Commons}.
\item Russell, S. (2019). \textit{Human Compatible}.
\item Amodei et al. (2016). \textit{Concrete Problems in AI Safety}.
\end{itemize}

\end{document}